\begin{table*}[pos=!tbp]\centering
\caption{Frozen solver configurations. All four internal variants share the same hierarchy and relaxation implementation.}\label{tab:configuration}
\begin{tabular}{llll}
\toprule
Method & Fine labels / aggregation & Fine relaxation & Coarsest limit\\
\midrule
ML-CPMG & Persistent watershed, $B_{\max}=256$ & Pillar line / point &512\\
Min-threshold & Connected foreground $k>\min k$ & Same as ML-CPMG &512\\
Single-median & Connected foreground $\log k>\operatorname{med}(\log k)$ & Same as ML-CPMG &512\\
Algebraic-only & Heavy-edge aggregation from level 0 & Same as ML-CPMG &512\\
SA & Symmetric strength, $\theta=0.04$ & Symmetric Gauss--Seidel &10\\
RS & Classical strength, $\theta=0.25$ & Symmetric Gauss--Seidel &512\\
\bottomrule
\end{tabular}
\par\smallskip\footnotesize
Internal methods use $\omega_P=0.45$, $\tau_P=0$ on the fine level and $0.04$
thereafter, full positive-edge lumping, $\eta_{\rm stall}=0.33$, one pre- and
one post-smoothing step, and at most eight nodes per algebraic aggregate.
The automatic rule selects line relaxation when the largest directional median
coupling is at least four times the next largest. On SPE10, calibration uses
the automatic rule and validation prescribes vertical lines. On OLYMPUS\_1 all
four internal methods use vertical lines. The three geometric methods also
semi-coarsen along pillars. Persistence, barrier, background and pedestal
parameters are $16$, $8$, $0.22$ and $2$, respectively. SA thresholds $0$ and
$0.1$ are reported in the sensitivity and faulted-reservoir studies.\normalsize
\end{table*}
